\begin{abstract}
For an elementary rule $f$, the second-order (Fredkin) automaton $F_f$ on $(\{0,1\}^2)^{\Z}$ is reversible, and the automorphism group of its spacetime subshift is the centralizer of $F_f$ in the automorphism group of the full $4$-shift. Many of these automata have walls, patterns that stay fixed or periodic at a fixed position whatever surrounds them; they were listed by Takesue (1989), who observed that no information passes through them. We determine what walls imply for the automorphism group. If a reversible automaton of radius $1$ has a word that it maps to itself in every context, and at least two symbols can sit between two copies of it, then permuting the contents of trains of identical segments embeds $\bigoplus_mS_m$, and hence every finite group, in the automorphism group. With Takesue's fixed walls this applies to $32$ of the $88$ conjugacy classes of second-order elementary automata, including $54$R and $37$R, and we prove that no other class has a fixed word of any length. Variants of the construction, for walls that are periodic in time or that carry a bit invisible to their neighbours, give the same conclusion in seven further classes. Three affine classes, among them that of $90$R, which has no wall and no blocking word, reach it through marker permutations; altogether $42$ classes ($102$ rules) have every finite group in their automorphism group. For four other affine classes the isomorphism rigidity theorem of Kitchens and Schmidt shows that every automorphism is affine, and the automorphism group is $\Z^2$, $\Z^2\times\Z/2$ or, for rule $150$, $\Z^3$. So for affine rules the absence of a blocking word does not force the automorphism group to be virtually $\Z^2$.
\end{abstract}
